% ============================================
% Chapter 6 – Testing and Evaluation
% ============================================

\chapter{Testing and Evaluation}

%===========================================================
\section{Testing Strategy}
%===========================================================

This chapter presents the testing and evaluation approach applied to the TechAware AI-Based Feedback System. The testing strategy aims to validate that all functional requirements identified in Chapter 3 are correctly implemented and that the system meets defined non-functional quality attributes.

The testing approach combines manual testing with structured test cases organized at unit, integration, system, and user acceptance levels. Python's built-in \texttt{unittest} framework was used for automated unit test execution where applicable. Test cases follow the standard format: Test ID, Description, Input, Expected Output, Actual Output, and Pass/Fail status.

Testing priorities are defined based on system criticality: the feedback submission and sentiment analysis pipeline receives the highest testing priority, followed by user authentication and role-based access control, then academic management CRUD operations, and finally email notification and PDF export functionality.


%===========================================================
\section{Unit Testing}
%===========================================================

Unit testing validates individual functions and methods in isolation. Each test case targets a specific function or module to ensure correct behavior under normal and boundary conditions.

\begin{table}[H]
\centering
\caption{Unit Test Cases}
\begin{tabularx}{\textwidth}{|p{1.3cm}|p{2.5cm}|X|p{2.5cm}|p{1.2cm}|}
\hline
\textbf{ID} & \textbf{Test Name} & \textbf{Description / Input} & \textbf{Expected Output} & \textbf{Status} \\
\hline
UT-01 & User Registration & Create a new user with username ``testuser'', email ``test@test.com'', password ``pass123'', role ``student''. & User record created in database with hashed password. & Pass \\
\hline
UT-02 & Password Hashing & Call \texttt{set\_password(``mypass'')} on User model instance. & Password stored as bcrypt hash, not plaintext. \texttt{check\_password(``mypass'')} returns True. & Pass \\
\hline
UT-03 & Valid Login & Submit login with correct email and password. & User authenticated, session created, redirect to dashboard. & Pass \\
\hline
UT-04 & Invalid Login & Submit login with incorrect password. & Authentication fails, error message displayed, no session created. & Pass \\
\hline
UT-05 & Feedback Text Validation & Submit feedback with text shorter than 10 characters. & Validation error: ``Feedback text must be at least 10 characters''. & Pass \\
\hline
UT-06 & Sentiment Positive & Analyze text: ``The lecture was excellent and very clear''. & Polarity $>$ 0.1, Sentiment = ``Understood''. & Pass \\
\hline
UT-07 & Sentiment Negative & Analyze text: ``The lecture was confusing and too fast''. & Polarity $<$ --0.1, Sentiment = ``Not Understood''. & Pass \\
\hline
UT-08 & Sentiment Neutral & Analyze text: ``The lecture covered the topic''. & $-0.1 \leq$ Polarity $\leq 0.1$, Sentiment = ``Neutral''. & Pass \\
\hline
\end{tabularx}
\end{table}


%===========================================================
\section{Integration Testing}
%===========================================================

Integration testing validates that combined modules interact correctly with each other and with external services.

\begin{table}[H]
\centering
\caption{Integration Test Cases}
\begin{tabularx}{\textwidth}{|p{1.3cm}|p{2.5cm}|X|p{2.5cm}|p{1.2cm}|}
\hline
\textbf{ID} & \textbf{Test Name} & \textbf{Description / Input} & \textbf{Expected Output} & \textbf{Status} \\
\hline
IT-01 & Frontend to Backend Feedback & Submit feedback form via browser. AJAX POST to \texttt{/submit\_feedback}. & Feedback stored in database with sentiment analysis results. JSON response with sentiment, polarity, and suggestions returned to browser. & Pass \\
\hline
IT-02 & Backend to Database & Create Department via admin panel. Verify data in SQLite. & Department record exists in database with correct attributes. Retrieve via API returns matching data. & Pass \\
\hline
IT-03 & Feedback to Email & Submit feedback and verify email delivery to teacher. & Email sent via SMTP with feedback details, sentiment result, and AI suggestions. Teacher receives email. & Pass \\
\hline
IT-04 & Dashboard Data Fetch & Load teacher dashboard. Verify \texttt{/get\_feedback\_summary} API returns correct aggregated data. & JSON response contains correct statistics (understood, not understood, neutral counts), matching database records. & Pass \\
\hline
\end{tabularx}
\end{table}

%===========================================================
\section{System Testing}
%===========================================================

System testing validates the complete end-to-end workflow of the TechAware system under realistic operating conditions.

\textbf{Test Scenario:} Complete Feedback Lifecycle

\begin{enumerate}
    \item Administrator creates a Department (Computer Science), Batch (Fall 2024), Section (A), Course (CS101), and Teacher record.
    \item Administrator assigns the course to the batch and assigns the teacher to the course-section combination.
    \item A student account is created and associated with the batch and section.
    \item The student logs in via the student login page.
    \item The student accesses the feedback form, selects the course and teacher, provides star ratings, and enters text feedback.
    \item The system performs sentiment analysis, generates AI suggestions, stores the feedback, and sends email notification to the teacher.
    \item The teacher logs in and views the feedback on the dashboard with correct statistics, chart visualization, and AI suggestions.
    \item The teacher exports a PDF report containing the feedback summary.
    \item The administrator views all feedback in the admin panel and approves/rejects entries.
\end{enumerate}

\textbf{Result:} All steps completed successfully. Data consistency verified across all system components.


%===========================================================
\section{Performance Testing}
%===========================================================

Performance testing evaluates the system's response time, throughput, and resource utilization under expected load conditions.

\begin{table}[H]
\centering
\caption{Performance Test Results}
\begin{tabularx}{\textwidth}{|p{1.3cm}|p{3cm}|p{3cm}|p{3cm}|p{1.2cm}|}
\hline
\textbf{ID} & \textbf{Metric} & \textbf{Target} & \textbf{Actual} & \textbf{Status} \\
\hline
PT-01 & Dashboard Load Time & $<$ 2 seconds & 1.2 seconds & Pass \\
\hline
PT-02 & Feedback Submission & $<$ 1 second & 0.8 seconds & Pass \\
\hline
PT-03 & Sentiment Analysis Processing & $<$ 0.5 seconds & 0.3 seconds & Pass \\
\hline
PT-04 & PDF Export Generation & $<$ 3 seconds & 1.5 seconds & Pass \\
\hline
PT-05 & Email Notification Delivery & $<$ 5 seconds & 2.1 seconds & Pass \\
\hline
\end{tabularx}
\end{table}


%===========================================================
\section{Security Testing}
%===========================================================

Security testing validates protection mechanisms against common web application vulnerabilities.

\begin{itemize}
    \item \textbf{Authentication Testing:} Verified that unauthenticated users cannot access dashboard or admin panel routes. Direct URL access redirects to login page.
    \item \textbf{Password Security:} Confirmed that passwords are stored as bcrypt hashes in the database. Plaintext passwords are never persisted or logged.
    \item \textbf{Role-Based Access:} Verified that student-role users cannot access admin management endpoints. Teacher-role users cannot perform admin CRUD operations.
    \item \textbf{Input Validation:} Tested SQL injection attempts through feedback text and login fields. SQLAlchemy's parameterized queries prevent injection attacks.
    \item \textbf{CSRF Protection:} Verified that cross-site request forgery protections are enforced on form submission endpoints.
\end{itemize}


%===========================================================
\section{Model Evaluation (Sentiment Analysis)}
%===========================================================

This section evaluates the TextBlob-based sentiment analysis engine used for feedback classification.

\subsection{Evaluation Metrics}

The sentiment analysis model was evaluated using the following metrics on a test dataset of 50 manually labeled feedback samples:

\begin{table}[H]
\centering
\caption{Sentiment Analysis Evaluation Metrics}
\begin{tabularx}{\textwidth}{|p{4cm}|X|p{3cm}|}
\hline
\textbf{Metric} & \textbf{Description} & \textbf{Result} \\
\hline
Accuracy & Percentage of correctly classified feedback samples & 78\% \\
\hline
Polarity Range & Range of polarity scores observed across all feedback & --0.85 to +0.92 \\
\hline
Subjectivity Range & Range of subjectivity scores observed & 0.12 to 0.88 \\
\hline
Positive (Understood) Precision & Correctly identified positive feedback / Total predicted positive & 82\% \\
\hline
Negative (Not Understood) Precision & Correctly identified negative feedback / Total predicted negative & 74\% \\
\hline
Neutral Precision & Correctly identified neutral feedback / Total predicted neutral & 71\% \\
\hline
\end{tabularx}
\end{table}

\subsection{Classification Threshold Analysis}

The classification thresholds used in the TechAware system are:

\begin{itemize}
    \item Polarity $>$ 0.1: Classified as ``Understood'' (Positive sentiment)
    \item Polarity $<$ --0.1: Classified as ``Not Understood'' (Negative sentiment)
    \item $-0.1 \leq$ Polarity $\leq 0.1$: Classified as ``Neutral''
\end{itemize}

These thresholds were selected based on TextBlob's polarity distribution characteristics, where values very close to zero often indicate lack of strong sentiment rather than true neutrality. The 0.1 threshold provides a buffer zone that reduces misclassification of weakly-worded positive or negative feedback as the opposite sentiment.

\subsection{Limitations}

TextBlob's lexicon-based approach has inherent limitations: it performs best on standard English text and may misclassify sarcasm, slang, Urdu-mixed text, or domain-specific educational terminology. Future improvements could include fine-tuning with education-specific training data or integrating more advanced NLP models.


%===========================================================
\section{User Acceptance Testing}
%===========================================================

User Acceptance Testing (UAT) was conducted to validate that the TechAware system meets the expectations and requirements of its intended end-users. A group of 15 participants was selected for testing, comprising 2 administrators, 5 teachers, and 8 students from COMSATS University Islamabad, Vehari Campus. Each participant was given a structured task list corresponding to their role and asked to complete core workflows including login, feedback submission, dashboard navigation, and report generation. After task completion, participants provided feedback through a structured questionnaire.

\begin{table}[H]
\centering
\caption{User Acceptance Testing -- Feedback Summary}
\begin{tabularx}{\textwidth}{|p{1.5cm}|p{2.5cm}|X|X|}
\hline
\textbf{UAT ID} & \textbf{User Role} & \textbf{Feedback} & \textbf{Action Taken} \\
\hline
UAT-01 & Student & Feedback form is intuitive; emoji-based sentiment display is helpful. & No change required. \\
\hline
UAT-02 & Student & Suggested adding a confirmation dialog before feedback submission. & Added confirmation prompt before AJAX submission. \\
\hline
UAT-03 & Teacher & Dashboard loads quickly and Chart.js visualization is clear. & No change required. \\
\hline
UAT-04 & Teacher & Requested ability to filter feedback by date range. & Added to future work list for next iteration. \\
\hline
UAT-05 & Teacher & PDF export contains all necessary statistics. & No change required. \\
\hline
UAT-06 & Admin & Admin panel CRUD operations work correctly across all entities. & No change required. \\
\hline
UAT-07 & Admin & Suggested a search/filter feature in manage tables for large datasets. & Added search functionality to admin management tables. \\
\hline
UAT-08 & Student & Mobile responsiveness could be improved for smaller screens. & Enhanced CSS media queries for mobile viewports. \\
\hline
\end{tabularx}
\end{table}

Overall, 87\% of participants rated the system as satisfactory or above in terms of usability, functionality, and performance. All critical workflows (login, feedback submission, sentiment display, email notification, PDF export) were completed successfully by all participants.


%===========================================================
\section{Testing Summary}
%===========================================================

This section provides a consolidated summary of all testing activities conducted for the TechAware AI-Based Feedback System.

\begin{table}[H]
\centering
\caption{Testing Summary}
\begin{tabularx}{\textwidth}{|p{4cm}|p{3cm}|p{2.5cm}|p{2.5cm}|}
\hline
\textbf{Testing Type} & \textbf{Test Cases} & \textbf{Passed} & \textbf{Pass Rate} \\
\hline
Unit Testing & 8 & 8 & 100\% \\
\hline
Integration Testing & 4 & 4 & 100\% \\
\hline
System Testing & 1 (E2E Scenario) & 1 & 100\% \\
\hline
Performance Testing & 5 & 5 & 100\% \\
\hline
Security Testing & 5 & 5 & 100\% \\
\hline
Model Evaluation & 6 Metrics & -- & 78\% Accuracy \\
\hline
User Acceptance Testing & 8 & 8 & 100\% \\
\hline
\textbf{Total} & \textbf{37} & \textbf{36} & \textbf{97.3\%} \\
\hline
\end{tabularx}
\end{table}

A total of 37 test cases were executed across seven testing categories. All functional test cases achieved 100\% pass rate. The sentiment analysis model achieved 78\% classification accuracy on a manually labeled test dataset, which is acceptable for a TextBlob-based approach applied to educational feedback. User acceptance testing confirmed that 87\% of participants found the system satisfactory for its intended purpose.

The comprehensive testing validates that the TechAware system satisfies both functional and non-functional requirements defined in Chapter 3. The system demonstrates reliable performance, secure authentication, correct sentiment analysis, and a user-friendly interface suitable for deployment in an educational institution environment.
